\section{Figures.sty}\label{sec:ref-figures}

TikZ, pgfplots, and one shared vocabulary of plot styles, so that every
drawn graph in every course looks like its neighbours. \textbf{No class
loads it.} pgfplots is the slowest thing in the toolkit, so a document
that draws says so, and one that does not pays nothing:

\begin{verbatim}
\usepackage{Figures}
\end{verbatim}

\subsection*{A figure is a file}

One \texttt{tikzpicture} per file, in the topic's \texttt{figures/}
directory -- a sibling of \texttt{problems/} and \texttt{images/} -- pulled
in with a plain \cs{import}:

\begin{verbatim}
\begin{center}
  \import{../figures/}{apex-figure-1-4-3.tex}
\end{center}
\end{verbatim}

\texttt{figures/} is for \LaTeX{} that draws something; \texttt{images/}
is raster. \cs{import} resolves relative to the consuming file, so the same
figure works from a lecture and from a problem file. As with
\cs{importproblem}, write the path literally, never through a macro.

A figure takes pgfplots' default size. To set several side by side, narrow
them for that one group:

\begin{verbatim}
\begin{center}
  % local to this group: center is one
  \pgfplotsset{width=0.36\linewidth}
  \import{../figures/}{a.tex}\hfill
  \import{../figures/}{b.tex}\hfill
  \import{../figures/}{c.tex}
\end{center}
\end{verbatim}

\subsection*{Where the vocabulary came from}

The style names are APEX Calculus's, from the figure preamble in that
book's source. Every figure in APEX is pgfplots code written against these
names, so with the package loaded \emph{an APEX figure compiles verbatim}:
copy the code out of the book's \texttt{<latex-image>} element, undo the
XML escaping (\texttt{\&lt;} and \texttt{\&gt;}), and it builds. Figures
drawn from scratch use the same names. APEX is CC~BY-NC~4.0; the package
carries the attribution, and a figure file whose code is the book's carries
a source line.

Tested against the 175 body figures of APEX chapters 1--4 and
sections~5.1 and~6.7. The handful that fail want something the package
deliberately does not load -- see \emph{Not loaded} below.

\subsection*{Curves, lines and dots}

\begin{center}
\begin{tabularx}{\linewidth}{@{}l Y@{}}
\texttt{firstcurvestyle}, \texttt{second\ldots}, \texttt{third\ldots} &
  the first, second and third curve on an axis: solid, dash-dotted,
  dash-dot-dotted, in the three role colours. A bare \cs{addplot} or
  \cs{addplot+} takes them in order from the cycle list, so most figures
  never name one.\\
\texttt{areastyle} & a half-transparent filled region.\\
\texttt{tangentline}, \texttt{normalline}, \texttt{secantline} &
  a line about a curve, with arrow tips; add \texttt{seg} to the name for
  the same line as a plain segment.\\
\texttt{asymptote} & a thick dashed line with arrow tips.\\
\texttt{guideline}, \texttt{symmetryline}, \texttt{lineseg} &
  thin black construction lines.\\
\texttt{soliddot}, \texttt{hollowdot} & marks only, not joined up: the
  function is defined here, or the curve heads here but is not here.\\
\end{tabularx}
\end{center}

\subsection*{Ends of a curve}

Added to an \cs{addplot} beside its curve style, to say what happens at
each end. A kite is ``it carries on''; a circle is an endpoint, filled if
it is included.

\begin{center}
\begin{tabularx}{\linewidth}{@{}l Y@{}}
\texttt{leftarrow}, \texttt{rightarrow}, \texttt{infinite} &
  a kite at the left end, the right end, or both.\\
\texttt{open}, \texttt{closed}, \texttt{openclosed}, \texttt{closedopen} &
  circles at both ends, named left end first.\\
\texttt{openleft}, \texttt{closedright}, \ldots & a circle at one end
  only.\\
\texttt{infiniteopen}, \texttt{closedinfinite}, \ldots & a kite at one end
  and a circle at the other.\\
\texttt{openinterval}, \texttt{closedinterval} & a heavy segment ending in
  parentheses or brackets, for a number line.\\
\end{tabularx}
\end{center}

\subsection*{Axes}

The one part of the package that changes pgfplots' behaviour unasked:
\emph{every} axis in a document that loads Figures gets the house default
-- axis lines through the origin rather than a box, labels $x$ and $y$,
one minor tick between labelled ones, and the curve cycle list. That is
what makes a figure file short: it states its window and its plots.
Override any of it per axis (\texttt{xlabel=\{\$t\$\}}).

\begin{center}
\begin{tabularx}{\linewidth}{@{}l Y@{}}
\texttt{numberline} & an axis with no $y$: a number line.\\
\texttt{xdiscontinuity}, \texttt{ydiscontinuity} & a break mark in an
  axis whose window does not start at~0.\\
\end{tabularx}
\end{center}

\subsection*{Knobs}

\begin{center}
\begin{tabularx}{\linewidth}{@{}l Y@{}}
\texttt{firstcolor}, \texttt{secondcolor}, \texttt{thirdcolor} & the three
  role colours (blue, red, a dark magenta). Every style is written in
  terms of them, so \cs{colorlet}\texttt{\{firstcolor\}\{black\}} after
  loading recolours every figure. APEX's greyscale print palette is in the
  package as three commented lines.\\
\texttt{compat=1.18} & pinned, not \texttt{newest}, so a \TeX{} Live
  upgrade cannot move every label; and set at all, because pgfplots warns
  on every build without it.\\
\end{tabularx}
\end{center}

\subsection*{Not loaded}

Only \texttt{arrows.meta} and \texttt{fillbetween} are loaded. APEX's own
preamble also loads \texttt{tikz-cd}, \texttt{tkz-euclide} and several
more TikZ libraries; each costs load time and a function graph needs none
of them. A figure that wants a \emph{library} loads it at the top of its
own file -- \cs{usetikzlibrary} works in the document body -- as the
handful of APEX figures that use \texttt{positioning} must. One that wants
a \emph{package} (\texttt{tkz-euclide}) needs it in the preamble of the
document that imports it.
